\section{Method}\label{sec:design}

The method separates three questions: what the requirements mean, what follows from that reading, and how to explain the result to the author. This lets us examine a translation error separately from a solver result or an inaccurate explanation. The detailed procedures appear in Appendix~\ref{app:procedures}; the main text follows the decisions that matter to the review.

\subsection{Pipeline overview}

\begin{figure}[H]
\centering
\resizebox{\linewidth}{!}{\begin{tikzpicture}[node distance=8mm]
\node[stage] (words) {Requirements\\\footnotesize words + rule labels};
\node[stage,right=of words] (formal) {Formal model\\\footnotesize symbols + formulas};
\node[stage,right=of formal] (finding) {Findings\\\footnotesize rules + witness};
\node[stage,right=of finding] (report) {Review report\\\footnotesize explanation + choices};
\draw[flow] (words)--node[above=6mm,annotation]{LLM + Z3}(formal);
\draw[flow] (formal)--node[above=6mm,annotation]{Z3}(finding);
\draw[flow] (finding)--node[above=6mm,annotation]{LLM}(report);
\node[annotation,below=4mm of words] {Input text};
\node[annotation,below=4mm of formal] {Saved model};
\node[annotation,below=4mm of finding] {Solver results};
\node[annotation,below=4mm of report] {Questions for the author};
\end{tikzpicture}
}
\caption{The three stages save their results separately. The solver stage makes no LLM calls.}\label{fig:pipeline}
\end{figure}

For authentication R3, the translation stage identifies the failed-attempt count and locked state. The solver finds that the lockout obligation cannot be obeyed together with the proposed administrator exception R5 at five failed attempts. The explanation stage returns the conflict and keep/change choices. Rule labels, formulas, assumptions and the witness remain available alongside the wording; the solver checks can be replayed without a model server.

\begin{table}[H]
\centering
\begin{tabularx}{\linewidth}{@{}lXX@{}}
\toprule
Responsibility & What it does & What it does not establish \\
\midrule
Language model & Extracts symbols and rules; proposes readings and assumptions; words findings & That the original author's intent was preserved \\
Python code & Chooses queries; clusters, votes, filters and combines results & A formal proof of the orchestration itself \\
Z3 & Returns satisfiability results, models and unsatisfiable cores for supplied formulas & Whether the symbols faithfully describe the real system \\
\bottomrule
\end{tabularx}
\caption{What each part of the system can establish.}
\end{table}

\subsection{Snapshot model}\label{sec:model}

Let $V$ be a finite vocabulary of Boolean, integer, real or finite-enumeration symbols. A valuation $\sigma$ assigns each symbol a value of its sort. It represents one snapshot of conditions, events and responses. There is no transition relation, trace, clock or built-in notion of \emph{previously}.

An atomic rule is $r_i=(a_i,c_i)$, with antecedent $a_i$ and consequent $c_i$. Its meaning is $\phi_i=a_i\Rightarrow c_i$; $\Phi=\bigwedge_i\phi_i$ denotes all retained rules. For the detailed solver examples, we now use the order case from Table~\ref{tab:order}. Its notation is local to that case: $s$ means submitted, $v$ means stock available, $f$ means fulfilled, $b$ means backordered, $c$ is a cancellation event, $q$ is the canceled-and-refunded state and $u$ means payments refunded. The standalone $c$ is distinct from the subscripted consequent $c_i$. The central rules are
\[
 \phi_1=(s\land v)\Rightarrow f,\qquad
 \phi_2=(s\land\neg v)\Rightarrow b,\qquad
 \phi_3=\True\Rightarrow\neg b.
\]
The complete five-rule model and its saved identifiers are in Appendix~\ref{app:formal-model}. The formulas use SMT-LIB syntax~\citep{smtlib} and are parsed into Z3 expressions; the checker does not execute arbitrary Python generated by the language model.

These definitions limit what a finding establishes. A conflict concerns the encoded obligations in an encoded situation. Whether the symbols capture the intended objects, timing and business rules remains a translation question.
\source{speccheck/logic.py; speccheck/checks.py; specs/order/formal.json}

\subsection{Translation and interpretation checks}\label{sec:formalize}

The first difficulty is choosing a vocabulary that preserves relationships between requirements. R2 and R3 must share the backorder symbol $b$; otherwise the contradiction becomes invisible. Extraction asks the model to preserve requirement IDs and source wording, and uses event markers to distinguish occurrences from states. Vocabulary validation checks names and sorts, but cannot verify that all requirements were found or that the abstraction is appropriate (Algorithm~\ref{alg:read}).

Repeated translation then asks whether the model gives a consistent reading within that vocabulary. The default is eight samples per rule at temperature 0.7. Z3 groups whole implications by equivalence. Here $\oplus$ is exclusive OR, and $\UNSAT$ means no satisfying assignment exists. The test proves that no assignment makes exactly one formula true; an unknown result proves neither equivalence nor inequivalence:
\[
 \phi\sim\psi\quad\Longleftrightarrow\quad\UNSAT(\phi\mathbin{\oplus}\psi).
\]
The largest group supplies the selected reading; the distribution over usable groups determines the initial accepted, ambiguous or rejected status. Rejected rules are reported for clarification and omitted from the main solver checks. Algorithms~\ref{alg:sample} and~\ref{alg:cluster} give the sampling, entropy thresholds and selection rules.

\example{A situation that distinguishes two readings.}
Let $n$ be the number of failed attempts. For ``lock the account when failed attempts reach five,'' $n=5$ and $n\ge5$ differ at six failed attempts with the account unlocked. That situation obeys the first implication and violates the second. In an illustrative eight-sample split, five readings of one kind and three of the other indicate disagreement; replacing $n=5$ with $5=n$ does not.

The pipeline uses such witnesses to ask focused questions about the original wording. It tests both presentation orders and retains answer counts, so the reviewer can see how a reading was challenged. Even unanimous preference for another reading is referred to the author; the implementation does not silently replace the requirement's meaning. Proposed alternatives still use the original vocabulary, so a missing meaning can remain undiscovered (Algorithm~\ref{alg:probe}).

An optional second model supplies another set of readings in the same vocabulary. Confirmation is conditioned on required event symbols, avoiding a spurious disagreement when a reader omits an explicit event guard. An unconfirmed accepted rule becomes ambiguous only when a usable alternative exists (Algorithm~\ref{alg:second-reader}). The second reader is off by default. Agreement within one model, agreement between models, and agreement across newly chosen vocabularies are different observations; none is a calibrated probability of correctness.
\source{speccheck/formalize.py; speccheck/probe.py}

\subsection{Relations between outcomes}

Boolean logic does not infer a relation from English names. Separate symbols for ``available'' and ``unavailable'' can both be true unless a formula relates them. The pipeline therefore asks for domain facts from symbol descriptions, without showing the requirements, and retains supported assumptions unless they are proved inconsistent with those already retained. These remain assumptions for the reviewer to assess (Algorithm~\ref{alg:facts}).

\example{Removal and retention.}
In the deletion case, let $d$ mean confirmation, $m$ removal and $k$ retention. The formulas $d\Rightarrow m$ and $k$ can coexist. If D2 is restated using the other outcomes while its private removal symbol is hidden, a supported reading $d\Rightarrow\neg k$ yields the assumed link
\[
 (d\land m)\Rightarrow\neg k.
\]
This can expose the clash with D4. It can also be wrong: ``remove'' may mean hiding a record while retaining an audit copy. The link procedure and subsequent probes preserve that qualification (Algorithms~\ref{alg:links} and~\ref{alg:probe}).

For rules whose triggering conditions are not proved incompatible, outcome-pair questions ask whether two definite outcomes, expressed with disjoint symbols, can coexist in one snapshot. The prompt does not include the triggering conditions. A majority of returned answers judging them incompatible adds the unguarded relation $\neg(c_i\land c_j)$ to the assumptions. The check skips disjoint events and already excluded pairs, and asks at most forty eligible pairs in document order. Solver unknown can leave a pair eligible; it does not prove compatible triggers. This mechanism is on by default in the frozen revision and is the subject of RQ1.
\source{speccheck/formalize.py; speccheck/probe.py}

\section{Findings}\label{sec:solver}

With an interpretation in place, the checker asks which obligations clash, which situations have no conditional response, and which allowed behaviours deserve review. The order example makes the distinction concrete: R2 and R3 conflict, an idle order may reasonably have no response, and fulfillment without submission may reveal a missing restriction.

\subsection{Conflicts}

Write $\SAT(F)$ for a satisfying solver result and $\UNSAT(F)$ for a proof that no satisfying assignment exists. Each query has a five-second timeout. An unknown result establishes neither. The search needs a feasible triggering situation as well as incompatible obligations.

For trigger group $G$, let $A_G=\bigwedge_{i\in G}a_i$. With an assumption set $B$, a conflict requires
\[
 \SAT\left(A_G\land\bigwedge B\right),\qquad
 \UNSAT\left(A_G\land\bigwedge B\land\Phi\right).
\]
For R2 and R3, $s\land\neg v$ is feasible, but it forces both $b$ and $\neg b$. Merely checking $\Phi$ would miss the problem: the solver can satisfy all five implications by choosing no submission or cancellation.

The search considers connected trigger groups of up to three rules and extracts a deletion-minimized unsatisfiable core only after an explicit unsatisfiable result. It excludes situations forbidden by one rule alone. The limit concerns trigger groups, not necessarily the number of rules in an explanation. A conclusive core is subset-minimal; it is not necessarily the smallest possible explanation, and the search does not enumerate every explanation (Algorithm~\ref{alg:conflicts}).

The checker runs first without assumptions, then with them. The second pass may add a conflict but cannot suppress one from the first; assumptions that exclude a rules-only trigger mark it doubtful. For example, assuming $\neg(f\land q)$ can make R1 alone forbid the submitted, in-stock, canceled-and-refunded situation. A rules-only pass preserves the R1/R5 clash for review instead of letting that assumption hide it.

\subsection{Gaps}
Let $C$ contain the rules with non-tautological antecedents. For an event $e$ to which a conditional rule reacts, the gap search seeks
\[
 e\land\bigwedge_{e'\ne e}\neg e'\land\bigwedge_{i\in C}\neg a_i.
\]
It returns at most one generalized example for each such event and one with no event. In change mode, a usable removed rule can seed the search with a situation it used to cover; that example takes priority for its selected event. At least one surviving conditional rule must still react to an event from the removed rule. Without reacted-to events, the fallback searches up to three generalized gaps, without removal attribution. If no conditional rule remains, it returns no gaps. Assumptions can mark a gap implausible but do not erase it; rejected translations can create apparent missing coverage (Algorithm~\ref{alg:gaps}).

In the order example, R1 and R2 cover submission with and without stock. The quiet case $s=c=q=\False$ has no conditional response, although R3 still forbids backordering; waiting may be intended. Removing R2 while retaining R1 leaves $s=\True,v=c=q=\False$ without a conditional response. This illustrates why event coverage is useful, while still being narrower than completeness of a business process.

\subsection{Scenarios}
For each rule, the allowed-scenario query is $\SAT(\Phi\land\neg a_i\land c_i)$. R1 says when fulfillment is required; it does not say ``only then.'' The model permits $s=c=q=b=\False$ with $f=\True$, so fulfillment without submission is a consequence for the author to review. A second search finds combined antecedents whose situation is forbidden by one rule alone. Both searches use the rules without domain assumptions. An LLM selects at most three consequences to discuss; offline replay without an LLM omits that presentation branch.
\source{speccheck/checks.py}

\subsection{Explanations and qualifications}\label{sec:explanation}

The explanation stage receives the finding, original requirements and an exact evidence sheet. It adds a title, explanation and keep/change choices. The witness and IDs remain available independently. An unparseable response leaves the finding unworded, while an API exception can still abort the run. The optional text review is off by default; its suggestions remain separate from solver findings.

The report also records sensitivity to alternative readings and vocabularies. It substitutes one ambiguous reading at a time and reruns the bounded checks; it does not test every combination. A gap annotation means that an alternative may cover part of the reported region, not that the whole gap disappears. Solver unknown can produce that hint, and a conflict absent from a bounded rerun does not prove consistency. Appendix~\ref{app:explanation} gives the counterexample and the rules for comparing multiple vocabularies.

All solver statements are conditional on the encoding. No proof certificate is archived, and generated prose is not formally checked. The R2/R3 result needs no assumed relation; a deletion conflict may depend on the removal/retention link; another finding may require simultaneous events. Reviewers need those distinctions alongside the explanation.
\source{speccheck/checks.py; speccheck/explain.py}

\section{Changes}\label{sec:change}

\subsection{Project baseline and proposed change}

OpenSpec keeps the project's existing behaviour under \code{openspec/specs/}. A proposed change has its own folder under \code{openspec/changes/}; its specification files describe the difference, or \emph{delta}, to apply to that baseline~\citep{openspec}. A capability such as authentication has the same relative path in both places. The checker loads the project's capability specifications, combines them with the selected change, and checks that combined document. It does not apply or archive the change in the source project.

The authentication example makes the layout concrete. These are the relevant files under \code{specs/auth_change/}:
\begin{lstlisting}
openspec/
|-- specs/auth/spec.md                     baseline
`-- changes/admin-never-locked/
    |-- proposal.md                       motivation
    `-- specs/auth/spec.md                delta
\end{lstlisting}

The \href{\authbaseurl}{baseline authentication specification on GitHub} contains the two requirements and four scenarios introduced in Table~\ref{tab:auth-baseline}. We return here to the authentication labels: R1--R4 are the existing scenarios in the combined full-mode document, and R5 is the proposed addition. These labels are separate from the order example's labels.

The \href{\authchangeurl}{proposed change on GitHub} adds one requirement. This is its literal specification text, with blank lines omitted:
\begin{lstlisting}
## ADDED Requirements
### Requirement: Administrator Access
Administrators MUST always be able to reach the system.
#### Scenario: Administrators are never locked out
- **WHEN** the account belongs to an administrator
- **THEN** the system does not lock the account
\end{lstlisting}

The four baseline scenarios remain \code{existing}; the administrator scenario becomes \code{added}, labelled R5 in this combined document. At five failed attempts, R3 requires locking an administrator's account under either threshold reading, while R5 forbids it. This particular conflict needs the base and the delta together. The saved findings also contain an older R1/R4 conflict: the checker retains it in the evidence but excludes it from the change report. The baseline is therefore not assumed to be free of defects.

\subsection{Applying the delta}

A delta file can use up to four operation sections. The checker applies them in the order shown in Table~\ref{tab:delta-operations}, regardless of their order in the file. Requirements are matched by name within their capability. A rename therefore takes effect before a removal or rewrite refers to the resulting name. References to missing requirements and additions of requirements already present are rejected (Algorithm~\ref{alg:change}).

\begin{table}[H]
\centering
\begin{tabularx}{\linewidth}{@{}llX@{}}
\toprule
Order & Section & What the delta supplies \\
\midrule
1 & \code{RENAMED} & A FROM/TO pair of requirement names; the requirement keeps its content under the new name. \\
2 & \code{REMOVED} & The name of an existing requirement. Text below that heading explains the reason for removal. \\
3 & \code{MODIFIED} & The full new requirement, including every scenario that should remain. \\
4 & \code{ADDED} & A new requirement's full text. ADDED can also create a specification for a capability with no baseline file. \\
\bottomrule
\end{tabularx}
\caption{The four operation headings have the form \texttt{\#\# ADDED Requirements}, with ADDED replaced by the appropriate operation.}\label{tab:delta-operations}
\end{table}

Within a MODIFIED requirement, scenario names identify which old and new scenarios correspond. The checker then compares their normalized bodies, ignoring case, extra whitespace and selected Markdown marks. Table~\ref{tab:scenario-status} gives the four possible states. These are textual classifications; they do not establish that a rewrite changes or preserves the intended meaning.

\begin{table}[H]
\centering
\begin{tabularx}{\linewidth}{@{}lX@{}}
\toprule
Scenario state & How it is assigned \\
\midrule
\code{existing} & The same name and normalized body remain, including inside a rewritten requirement. \\
\code{modified} & The name remains but the normalized body changes. \\
\code{added} & The scenario name is new, or its entire requirement is added. \\
\code{removed} & The whole requirement is removed, or its rewritten text omits an old scenario name. \\
\bottomrule
\end{tabularx}
\caption{Scenario provenance determines which findings are relevant to the change.}\label{tab:scenario-status}
\end{table}

For example, rewriting Account Lockout while retaining its Locked account scenario unchanged leaves that scenario \code{existing}. Changing only the Lockout threshold scenario makes that scenario \code{modified}. Omitting Locked account from the replacement text instead makes it \code{removed}. These are independent illustrations of the matching rules; the saved administrator change above only adds R5.

The default route formalizes the combined document, including removed scenarios labelled \code{X1}, \code{X2}, and so on. Before solver analysis, it separates the removed obligations and discards links or outcome-pair assumptions tagged with them. Ordinary voted facts remain. The removed obligations are no longer active rules, but their antecedents can guide a search for newly uncovered situations. Their text can still affect the full run's vocabulary and other assumptions.

\subsection{Removing a response}

\example{Removing the no-stock response.}
Take the order model and remove R2 while retaining R1. This is a constructed change to the order EARS example, not a saved OpenSpec order change. We use X1 for the removed R2 to explain the change-mode label. R1 still handles submission with stock. The situation $s=\True,v=c=q=\False$ now has no conditional response; R3 remains a prohibition, not a response. The removed rule's antecedent points directly to this missing branch.

\begin{figure}[H]
\centering
\resizebox{0.94\linewidth}{!}{\begin{tikzpicture}[node distance=9mm and 12mm]
\node[stage,text width=3.2cm] (active) {Surviving rules\\R1, R3, R4, R5};
\node[stage,text width=3.2cm,below=of active] (removed) {Removed X1 (old R2)\\No-stock response};
\node[stage,text width=3.8cm,right=of active] (check) {Active obligations\\checked as usual};
\node[stage,text width=3.8cm,right=of removed] (seed) {No conditional response\\$s=\True$\\$v=c=q=\False$};
\node[stage,text width=2.7cm,right=of seed] (finding) {Gap from X1\\\texttt{opened\_by}};
\draw[flow] (active) -- (check);
\draw[flow] (active.south east) -- (seed.north west);
\draw[flow] (removed) -- node[annotation,above] {antecedent} (seed);
\draw[flow] (seed) -- (finding);
\end{tikzpicture}
}
\caption{The removed no-stock rule leaves the active model, but supplies the situation used to investigate a gap. The labels illustrate the order change.}\label{fig:removal}
\end{figure}

\begin{lstlisting}
{"type": "gap", "event": "order_submitted",
 "opened_by": ["X1"], "about_change": true}
\end{lstlisting}
These fields illustrate the attribution; they are not a literal saved order-change result. The change report includes this gap even though X1 is no longer an active obligation. If every response to submission is removed, the event is treated as retired and this search does not report a removal gap. State-only removals and rejected removed translations also fall outside this attribution mechanism. A gap associated with an added or modified rule's event is relevant to the change; that association alone does not prove the change caused it.

\subsection{Saved-base reuse}\label{sec:incremental}

The optional \code{--incremental} route formalizes the existing specifications once and retains that base model. It then translates only added or rewritten scenarios, labelled \code{N1}, \code{N2}, and so on, using the base vocabulary. Existing symbols and readings stay fixed. The pipeline probes the new rules and their eligible interactions with old rules, combines the models, removes deleted or replaced rules, and reruns the ordinary solver analyses. It reuses translation work; it does not incrementally update solver proofs (Algorithm~\ref{alg:incremental}).

\begin{figure}[H]
\centering
\resizebox{\linewidth}{!}{\begin{tikzpicture}[node distance=9mm and 9mm]
\node[stage,text width=2.8cm] (full) {Full run (default)\\Whole merged text};
\node[stage,text width=3.3cm,right=of full] (translate) {Choose vocabulary\\translate all rules};
\node[stage,text width=3.3cm,right=of translate] (analyze) {Active model\\Run solver checks};
\node[stage,text width=2.8cm,below=of full] (base) {Optional reuse\\Saved base model};
\node[stage,text width=3.3cm,right=of base] (new) {Scenario edits\\Translate and probe\\Keep old readings};
\node[stage,text width=3.3cm,right=of new] (combine) {Combine models\\Drop removed or\\replaced rules};
\draw[flow] (full) -- (translate);
\draw[flow] (translate) -- (analyze);
\draw[flow] (base) -- (new);
\draw[flow] (new) -- (combine);
\draw[flow] (combine) -- (analyze);
\node[annotation,below=4mm of base,text width=3.5cm] {Build the base once; reuse while\\base text and model name match};
\end{tikzpicture}
}
\caption{Two routes to the same analysis procedures. Their resulting encodings may differ, so the same solver checks do not establish equal detection quality.}\label{fig:incremental}
\end{figure}

For the administrator exception, the saved model supplies the lockout reading and vocabulary; only the new exception and its interactions need new translation questions. Repeated alternatives against the same base can share that initial work. Editing or archiving into the base specifications changes its text fingerprint and requires another base formalization. The cache also checks the model name, but not sampling settings, prompts, endpoint or code version. Invalid labels or an unusable model abort; there is no automatic fallback to a full merged-document run.

This mode currently follows scenario edits. Changing only a requirement's SHALL sentence while retaining the same scenarios, or adding a requirement without scenarios, can yield no new text and reuse the old model. Full mode is needed for those changes and for reconsidering existing interpretations. Incremental mode also excludes multiple vocabularies, a second reader and text review. E20 measures the resulting detection and request tradeoff; it does not establish equivalent semantic accuracy.
\source{speccheck/change.py; speccheck/cli.py}
